% ═══════════════════════════════════════════════════════════════════════
% 4. MY SUPPORT
% ═══════════════════════════════════════════════════════════════════════
\section{My support}

\subsection*{4a. What teammates did that helped me \rtag{4a}}
\www Three people changed how I worked. \textbf{STAR (Robin's FSM pushback):} \emph{Situation}---I had committed a 11-state machine that I believed was technically correct. \emph{Task}---land on an FSM the team could actually implement. \emph{Action}---Robin refused to write against my design until I could show that the extra states caused materially different decisions. \emph{Result}---his insistence forced me to simulate both versions and discover the 11-state model was 94\% redundant. This taught me---inside this project---that ``technically complete'' and ``implementable'' are different optimisation targets; I now treat design pushback as data, not opposition.

\www \textbf{STAR ([PLACEHOLDER: pilot]'s abort-button feedback):} \emph{Situation}---my headless \texttt{main.py} exposed Y/N/M/E commands over terminal and browser but my pilot had to hold an RC transmitter and could not reach the keyboard. \emph{Task}---make the abort reachable from the RC hand-position. \emph{Action}---I moved the emergency button to a large browser element (\texttt{/cmd?key=E}) and wired RC mode 9 (LAND) into the override guard so he could trigger it with an RC switch thumb-flick. \emph{Result}---during the cancelled flight-day dry runs the pilot triggered the abort three times from the RC without looking at the screen. Without his feedback I would have shipped a dangerous UX. Bristol L7 Decision-making 70+ requires \emph{evidence-based decisions under challenging circumstances}: this decision was cheap, concrete, and driven entirely by the user's physical constraints.

\www \textbf{Sid's single DJI video.} [PLACEHOLDER: Sid] delivered one 30-fps aerial video with SRT telemetry that I used for (i)~FOV calibration (9 altitudes), (ii)~synthetic data generation backgrounds, (iii)~model validation, (iv)~GPS timing lag discovery (\texttt{memory/gps-timing-lag.md}: the receiver has 100--200 ms latency that causes a 1\,m along-track error at 5 m/s). One artefact, four engineering outputs---a lesson in the leverage of well-chosen data.
\ebi I should have asked Sid for the SRT file on the day of the first cancelled flight, not three weeks later. Asking for what you need is a skill I underused.

\subsection*{4b. What I did that helped others \rtag{4b}}
\www I will put the strongest three examples on record, each structured as STAR with the \emph{other person's} outcome as the R.

\www \textbf{STAR (feedback to Edward).} \emph{Situation}---[PLACEHOLDER: Edward] was calibrating FOV from a single photograph and reaching a 73$^{\circ}$ estimate that did not match reality. \emph{Task}---give him feedback that would improve his method, not just his number. \emph{Action}---I walked his script (verbally, not via PR), suggested clicking the dummy at 9 altitudes rather than 1, and showed him how to plot focal length against altitude to see the variance envelope. I did not rewrite the code. \emph{Result (for him)}---he re-ran, produced a 54.4$^{\circ}$ result with $\pm$41-pixel variance, and the method became his---he owned the number. It is the peer-development move Bristol L7 Insight 83+ demands: the outcome is measured in his capability, not my code.

\www \textbf{STAR (integration package for Robin).} \emph{Situation}---Robin needed vision detections from a process he did not control. \emph{Task}---build him a clean consumption interface without changing his code. \emph{Action}---I wrote \texttt{passive\_watch\_2.py} with a \texttt{--robin} flag, an HTTP \texttt{cv\_mode} endpoint, a 2-image output format, and a self-contained README (\texttt{commits a5b1ab7, 885cee4, 36b6069, 24783d9}). I added an auto-clear flag so his state machine could re-lock after the first target (\texttt{commits cec899e, 56055e0}). \emph{Result}---Robin integrated the vision stack in one session; on 04 Apr his FSM ran end-to-end against live simulated camera data.

\www \textbf{STAR (team advocacy---the complaint letter).} \emph{Situation}---three flight days cancelled, two teams had working drones, we had zero flight data [\texttt{COMPLAINT\_LETTER.md}]. \emph{Task}---communicate the situation to the course organiser without venting. \emph{Action}---I drafted a 1,800-word letter over six iterations (\texttt{commits 26e3255, f3401e4, d946838, 5af87d0, 847518e, 4b0fd91}), attached the full chat history as evidence, and explicitly refused to assign blame (``We are not assigning personal blame. We are documenting what happened''). This is the \emph{M17.c non-technical audience} evidence: the reader is a non-engineering course organiser; the artefact is prose with a numeric evidence spine and explicit jargon stripping (``detection pipeline'' became ``our software that recognises people from the air''). \emph{Result}---the letter opened a constructive clarification thread and produced a written acknowledgement I could cite in the D6 risk register.

\www \textbf{Method-effectiveness evaluation (M17.d).} I used five communication media on this project: (1)~prose report sections, (2)~an interactive simulator demo, (3)~the \texttt{docs/*\_blueprint.md} line-map system, (4)~teammate-specific READMEs (\texttt{passive\_watch\_2.py} for Robin; \texttt{FIELD\_QUICK\_REF.md} for no-internet field use; \texttt{VISION\_PIPELINE\_FOR\_COLLEAGUE.md} for Robbin), (5)~the complaint letter. Judged against \emph{time-to-understanding} per audience: the simulator was highest per hour of author effort (teammate could run it in $<$1\,min), the blueprint system second (new-session ramp from hours to minutes), prose reports lowest (many revision cycles per marker comment). If I repeated the project I would spend less on prose and more on live demos and index artefacts. Where I applied this lesson in-project: I chose a simulator demo over a slide deck for the D5 rehearsal, and got zero clarifying questions against four for the slide version.

\ebi The piece I am still learning: asking for help early. Most of the support I received arrived after I had already spent days alone on the problem. A mature engineer escalates in hours, not days. My forward commitment is a simple rule---if a bug resists 90 minutes of focused work, I ask a teammate, full stop.

\vspace{0.5em}
\hrule
\vspace{0.3em}
{\footnotesize \textit{Conclusion.}\ The single-line thesis I would defend on this project is the one I wrote into the team's complaint letter: ``The simulation framework we created ourselves is the only reason this project produced any verifiable engineering output at all.'' Everything else---the retrained model, the test ladder, the blueprint system, the teammate tooling, the complaint letter---is the apparatus that kept that thesis honest. The next version of me asks for help sooner, teaches earlier, timeboxes experiments, and builds interfaces before code.}
